% Companion Manuscript for Figure 3 of the Hypothesis Surface paper.
%
% Merged content from three working documents:
%   - obs_bound_scaling_proposition.tex (universal exponent theorem)
%   - bicomplex_carrier_v1.tex (minimal curved bicomplex setup only)
%   - test_time_cohomology.tex (observer-bounded cohomology reframe)
%
% Numbering is designed so the main text's patch references land
% correctly:
%   - "companion MS, Thm. 1" = Universal Critical-Exponent Theorem (§2)
%   - "companion MS, Def. 5" = Observer-Bounded Cohomology (§7)
%
% The bicomplex content is included only to the minimal extent needed
% to define observer-bounded cohomology as the associated graded of
% an Omega_O-filtered curved bicomplex. The full bicomplex construction
% with cacophony/FRW cocycle placement is Principia Book IX material
% (post-deadline).
%
% Proofs of Thm 2 and Thm 3 are abbreviated relative to the standalone
% working documents; full derivations remain in the source files above.

\documentclass[11pt]{article}
\usepackage[margin=1in]{geometry}
\usepackage{amsmath, amssymb, amsthm}
\usepackage{mathtools}
\usepackage{enumitem}
\usepackage{hyperref}

\newtheorem{definition}{Definition}
\newtheorem{theorem}{Theorem}
\newtheorem{proposition}{Proposition}
\newtheorem{lemma}{Lemma}
\newtheorem{corollary}{Corollary}
\newtheorem{remark}{Remark}
\newtheorem*{principle}{Principle}

\newcommand{\Obs}{\mathcal{O}}
\newcommand{\OmegaO}{\Omega_{\Obs}}
\newcommand{\kappaO}{\kappa_{\Obs}}
\newcommand{\AO}{A_{\Obs}}
\newcommand{\nablaO}{\nabla_{\Obs}}
\newcommand{\hO}{h_{\Obs}}
\newcommand{\gO}{g_{\Obs}}
\newcommand{\DO}{D_{\Obs}}
\newcommand{\rmin}{s_{\Obs}}
\newcommand{\EE}{E}
\newcommand{\Kt}{\mathcal{K}_t}
\newcommand{\End}{\mathrm{End}}
\newcommand{\Real}{\mathbb{R}}
\newcommand{\Comp}{\mathbb{C}}
\newcommand{\HO}{H_\Obs}
\newcommand{\ObsSys}{\mathsf{ObsSys}}
\newcommand{\Frob}{\mathsf{FRW}}
\newcommand{\Cacph}{\mathsf{Cac}}
\newcommand{\Sym}{\mathrm{Sym}}

\title{Companion Manuscript, Figure~3: \\
Observer-Bounded Access, the Universal Critical-Exponent
Theorem, and Observer-Bounded Cohomology\\[0.4em]
\large Proofs supporting Figure~3 and the Empirical Concurrence
paragraph of the Hypothesis Surface manuscript (\S7,
\texttt{main.tex}).}
\date{}

\begin{document}
\maketitle

\begin{abstract}
\noindent This companion manuscript provides the rigorous
statements and proofs supporting Figure~3 of the Hypothesis
Surface paper. The core result is the
\emph{Universal Critical-Exponent Theorem} (\S\ref{sec:universal},
Thm.~\ref{thm:universal}): the cacophony cliff's
$\delta_{\min} \sim \varepsilon^{-1/2}$ scaling and the FRW past
light cone's $+1/2, +2/3$ scaling are \emph{instances of the same
theorem}, recovered by evaluating a single functor
$\OmegaO:\ObsSys \to C^\infty(I)$ on the cacophony bilinear form
and the FRW spatial metric respectively. The matter-era $-2/3$
is a prediction, not a regime-distinct curiosity.
The second half of the manuscript reframes the obstruction to
strict cohomological unification as \emph{content}: the
Chern--Weil class that prevents one-shot decurving is
structurally the same object that expresses the observer bound,
so under observer-bounded cohomology
(\S\ref{sec:cohom}, Def.~\ref{def:ocoh}) the obstruction becomes
the natural carrier of the shared invariant rather than an
impediment to it. Preliminary definitions depend on the
bundle-valued formalism of the companion manuscript
\texttt{fig2\_companion\_fftc.tex}.
\end{abstract}

\section*{0. Scope, method, and honest framing}
\label{sec:scope}
\phantomsection\addcontentsline{toc}{section}{0. Scope, method, and honest framing}

\noindent\textbf{What this document is.}
A rigorous unification of two empirical $\sqrt{\cdot}$-scalings
under a single observer-bounded access functor, together with
the cohomological framework in which the resulting obstruction is
content rather than failure. Both halves support
Figure~3 of the Hypothesis Surface paper.

\noindent\textbf{What this document is not.}
A claim that the cacophony cliff and the FRW cosmology are
``the same phenomenon.'' They are formally distinct
observer-bounded systems that happen to be instances of the same
functorial construction, with explicitly different parameter
content. Uniqueness of the underlying mechanism is not asserted.

\noindent\textbf{Assumptions made explicit.}
(i) Cacophony uses the one-parameter family of Gram matrices
$\Gamma(\rho) = (1 - \rho) I + \rho J$ of~\cite{tiffany2026cacophony}.
(ii) FRW is perfect-fluid, flat spatial slice, $\Lambda$CDM-calibrated
by Planck~2018. (iii) Observer bounds are encoded by a
resolution-floor $\rmin > 0$ and an anti-masking rate $\mu$
controlling spectral faithfulness
(\texttt{main.tex}, Def.~Inv~2).

\noindent\textbf{What would falsify the unification.}
(a) An equation-of-state $w$ at which the universal exponent
$-2/(3(1+w))$ fails empirically, given Planck-calibrated ambient
geometry. (b) A cacophony-type bilinear form whose degenerating
exponent is not $p = 1$, or a rank-one structure whose critical
scaling is not $-1/2$. Both are measurable.

\section{Category of observer-bounded systems}
\label{sec:obsys}

\begin{definition}[Observer-bounded system]
\label{def:obsys}
An \emph{observer-bounded system} is a tuple
$(V, \langle\cdot,\cdot\rangle, I, B, W)$, where:
\begin{enumerate}[label=(\alph*),leftmargin=*,itemsep=2pt]
\item $V$ is a finite-dimensional real inner product space with
inner product $\langle\cdot,\cdot\rangle$.
\item $I \subset \Real$ is an interval (the \emph{control
parameter}), and $B: I \to \Sym^+(V)$ is a continuously
varying symmetric positive-semidefinite bilinear form on $V$
(the \emph{capacity form}).
\item $W \subseteq V$ is a nonzero linear subspace (the
\emph{observer-accessible subspace}).
\end{enumerate}
Morphisms $(V, \langle,\rangle, I, B, W) \to
(V', \langle,\rangle', I', B', W')$ are pairs $(\phi, \psi)$ where
$\phi: V \to V'$ is a linear isometry with $\phi(W) \subseteq W'$
and $\psi: I \to I'$ is a continuous map such that
$\phi^* B'_{\psi(\lambda)} = B_\lambda$ on $W$. Composition is
by pairwise composition. Let $\ObsSys$ denote the resulting
category.
\end{definition}

\begin{definition}[Observer-bounded access functor]
\label{def:accessor}
The \emph{observer-bounded access functor} is
\[
\OmegaO: \ObsSys \to C^\infty(I; \Real_{\geq 0}),
\qquad
\OmegaO(V, \langle,\rangle, I, B, W)(\lambda) \;:=\;
\inf_{\substack{w \in W \\ \|w\| = 1}}
\langle w, B_\lambda w \rangle,
\]
i.e., the smallest Rayleigh quotient of $B_\lambda$ restricted
to the observer-accessible subspace $W$. $\OmegaO$ sends
morphisms by pullback: $\OmegaO(\phi, \psi)(f)(\lambda) :=
f(\psi(\lambda))$.
\end{definition}

\begin{lemma}[Rayleigh dual]
\label{lem:rayleigh}
$\OmegaO$ is well-defined, functorial, and satisfies the
\emph{Rayleigh dual}: the observer-bounded access cost
$\kappaO^{\mathrm{access}}(\lambda) :=
\sup_{w \in W, \|w\|=1} \langle w, B_\lambda w\rangle^{-1/2}$
satisfies
\[
\kappaO^{\mathrm{access}}(\lambda) \;=\; \OmegaO(\cdot)(\lambda)^{-1/2}.
\]
Access cost diverges exactly when capacity vanishes.
\end{lemma}

\begin{proof}
The infimum exists because $B_\lambda$ is continuous and the unit
sphere in $W$ is compact. Functoriality is immediate from the
definition of morphisms. The Rayleigh dual is the minimax
identity: for symmetric positive semidefinite $B_\lambda$ and
$w \in W$, $\langle w, B_\lambda w\rangle^{-1/2}$ is minimized
over unit $w$ precisely where $\langle w, B_\lambda w\rangle$ is
maximized, and \emph{diverges} precisely where
$\langle w, B_\lambda w\rangle$ attains zero, which is the
infimum definition of $\OmegaO$.
\end{proof}

\section{Universal critical-exponent theorem}
\label{sec:universal}

\begin{theorem}[Universal Critical-Exponent]
\label{thm:universal}
Let $\mathcal{S} = (V, \langle,\rangle, I, B, W)$ be an
observer-bounded system, and suppose there exists
$\lambda^* \in \overline{I}$ such that
$\OmegaO(\mathcal{S})(\lambda) \to 0$ as $\lambda \to \lambda^*$
with
\[
\OmegaO(\mathcal{S})(\lambda) \;=\; c\,|\lambda - \lambda^*|^p \;+\;
o\!\left(|\lambda - \lambda^*|^p\right),
\qquad c > 0,\; p > 0.
\]
Then the access cost $\kappaO^{\mathrm{access}}(\lambda)$ diverges
at $\lambda^*$ with the \emph{universal exponent} $-p/2$:
\[
\kappaO^{\mathrm{access}}(\lambda) \;=\; c^{-1/2}\,
|\lambda - \lambda^*|^{-p/2} \;+\;
o\!\left(|\lambda - \lambda^*|^{-p/2}\right).
\]
In particular, the critical exponent depends only on the
vanishing order $p$ of $\OmegaO$, not on the ambient geometry or
the particular form of $B_\lambda$.
\end{theorem}

\begin{proof}
Immediate from Lemma~\ref{lem:rayleigh}: if
$\OmegaO(\mathcal{S})(\lambda) = c\,|\lambda - \lambda^*|^p (1 + o(1))$
then
$\OmegaO(\mathcal{S})(\lambda)^{-1/2} =
c^{-1/2} |\lambda - \lambda^*|^{-p/2}(1 + o(1))^{-1/2}$,
and the right-hand side has the same leading-order exponent
$-p/2$ because $(1 + o(1))^{-1/2} = 1 + o(1)$.
\end{proof}

\begin{remark}[The $1/2$]
The universal exponent is $-p/2$ rather than $-p$ or $-1/p$
because \emph{access cost is the square root of inverse
capacity}. This is a consequence of the Rayleigh dual and is
independent of any physical interpretation. The $1/2$ appears
wherever capacity enters quadratically into an access-resistance
functional.
\end{remark}

\begin{corollary}[Linear vanishing gives exponent $-1/2$]
\label{cor:linear}
If $\OmegaO$ vanishes linearly ($p = 1$) at $\lambda^*$, then
$\kappaO^{\mathrm{access}}(\lambda) \sim |\lambda - \lambda^*|^{-1/2}$.
\end{corollary}

\begin{corollary}[Monomial vanishing]
\label{cor:monomial}
If $B_\lambda$ has a simple zero eigenvalue along a one-dimensional
direction in $W$ at $\lambda^*$ with eigenvalue vanishing as
$|\lambda - \lambda^*|^q$, then $\OmegaO \sim |\lambda - \lambda^*|^q$
(so $p = q$) and the exponent is $-q/2$.
\end{corollary}

\section{Cacophony as instance ($p = 1$)}
\label{sec:cacophony}

\begin{proposition}[Cacophony as observer-bounded system]
\label{prop:cac}
The multi-constraint LLM alignment problem of~\cite{tiffany2026cacophony}
is an observer-bounded system $\Cacph(k)$ as follows:
\begin{enumerate}[label=(\roman*),leftmargin=*,itemsep=2pt]
\item $V = \Real^k$ with the Euclidean inner product;
\item $I = [0, 1/(k-1))$ and the control parameter is $\rho \in I$;
\item $B_\rho = (1+\rho) I_k - \rho J_k$ (equivalently,
$I_k - \rho(J_k - I_k)$), where $J_k$ is the all-ones matrix;
this is the rank-one perturbation whose smallest eigenvalue
governs the cacophony feasibility cliff of~\cite{tiffany2026cacophony};
\item $W = V = \Real^k$ (all coordinates are observer-accessible).
\end{enumerate}
Then $\OmegaO(\Cacph(k))(\rho) = 1 - \rho(k-1)$, which vanishes
linearly ($p = 1$) at $\rho^* = 1/(k-1)$.
\end{proposition}

\begin{proof}
$B_\rho$ has eigenvalue $1 - \rho(k-1)$ on the all-ones vector
(since $(1+\rho)\mathbf{1} - \rho J_k \mathbf{1} = (1+\rho-\rho k)
\mathbf{1} = (1-\rho(k-1))\mathbf{1}$), and eigenvalue $1+\rho$
on its orthogonal complement (multiplicity $k-1$, since
$J_k v = 0$ for $v \perp \mathbf{1}$). For $\rho \in [0, 1/(k-1))$,
the smallest eigenvalue over $W = V$ is $1 - \rho(k-1)$, attained
on the all-ones direction, vanishing linearly at
$\rho^* = 1/(k-1)^-$. Thus $\OmegaO(\Cacph(k))(\rho) = 1 - \rho(k-1)$
and $p = 1$.
\end{proof}

\begin{corollary}[Cacophony cliff]
\label{cor:cliff}
$\kappaO^{\mathrm{access}}(\Cacph(k), \rho)
\sim (1/(k-1) - \rho)^{-1/2}$ as $\rho \to 1/(k-1)^-$. This is
the cacophony cliff.
\end{corollary}

\begin{proof}
Theorem~\ref{thm:universal} with $p = 1$.
\end{proof}

\section{FRW cosmology as instance ($p = 4/(3(1+w))$)}
\label{sec:frw}

\begin{proposition}[FRW perfect-fluid cosmology as observer-bounded system]
\label{prop:frw}
The past light cone of an observer in a flat FRW cosmology with
perfect-fluid equation of state $w$ (so $p(t)/\rho(t) = w$,
constant) is an observer-bounded system $\Frob(w)$ as follows.
Let $a(t)$ solve the Friedmann equations with initial condition
$a(t_0) = 1$; the scale factor satisfies
$a(t) \propto t^{2/(3(1+w))}$ for $w \neq -1$.
\begin{enumerate}[label=(\roman*),leftmargin=*,itemsep=2pt]
\item $V = T_x M^{(3)}$, the tangent space of the observer's
spatial slice at its worldline;
\item $I = (0, t_0]$ (cosmic time, $t_0$ = present);
\item $B_t = a(t)^2 \gamma$, where $\gamma$ is the flat spatial
metric on the slice;
\item $W = V$, the observer's spatial tangent space. (The
past-directed tangent cone at the observer spans $V$; since
$B_t$ is homothetic, the Rayleigh quotient is constant over
directions, so replacing the cone by its linear span $V$ leaves
$\OmegaO$ unchanged while satisfying the linear-subspace
requirement of Def.~\ref{def:obsys}(c).)
\end{enumerate}
Then
\[
\OmegaO(\Frob(w))(t) \;=\; a(t)^2 \;\propto\; t^{4/(3(1+w))},
\]
which vanishes at $t^* = 0$ (the Big Bang) with vanishing order
$p = 4/(3(1+w))$.
\end{proposition}

\begin{proof}
The Rayleigh quotient of $B_t = a(t)^2 \gamma$ over unit vectors
in $W$ (with $\|\cdot\|$ the inner-product norm on $V$,
normalized by $\gamma$ so that $a(t_0)^2 = 1$) is constant over
directions because the metric is homothetic with conformal factor
$a(t)^2$: $\langle w, B_t w\rangle = a(t)^2 \gamma(w, w) = a(t)^2
\|w\|^2 = a(t)^2$. Thus $\OmegaO(\Frob(w))(t) = a(t)^2$. For a
flat perfect-fluid FRW universe, $a(t) \propto t^{2/(3(1+w))}$,
so $a(t)^2 \propto t^{4/(3(1+w))}$, vanishing at $t^* = 0$ with
order $p = 4/(3(1+w))$.
\end{proof}

\begin{corollary}[Radiation-era instance]
\label{cor:rad}
At $w = 1/3$ (radiation), $p = 4/(3(1 + 1/3)) = 1$. So
$\kappaO^{\mathrm{access}}(\Frob(1/3), t) \sim t^{-1/2}$ and the
access-cost exponent is $-1/2$.
\end{corollary}

\begin{corollary}[Matter-era instance]
\label{cor:mat}
At $w = 0$ (matter), $p = 4/3$. So
$\kappaO^{\mathrm{access}}(\Frob(0), t) \sim t^{-2/3}$ and the
access-cost exponent is $-2/3$.
\end{corollary}

\begin{corollary}[General perfect-fluid equation-of-state scaling]
\label{cor:genw}
For arbitrary $w \neq -1$, the universal exponent evaluates to
$-p/2 = -2/(3(1+w))$. In particular, this recovers the
radiation $-1/2$, matter $-2/3$, and extrapolates smoothly to
stiff fluids ($w \to 1$: exponent $\to -1/3$) and cosmological
constant ($w \to -1$: exponent $\to -\infty$; the expansion is
non-power-law and the asymptotic form changes, so the
theorem doesn't apply to $w = -1$ without reformulation).
\end{corollary}

\begin{remark}[Matter-era $2/3$ as prediction]
\label{rem:matter-prediction}
The figure caption in \texttt{main.tex} previously framed the
matter-era $2/3$ as ``regime-distinct'' from the cacophony
cliff's $1/2$. Under Theorem~\ref{thm:universal},
both exponents are instances of the same functor with different
vanishing orders $p$. The matter-era $2/3$ is a
\emph{prediction} of the universal theorem given $w = 0$, not a
separate regime. The framework gains predictive content across
the full equation-of-state range.
\end{remark}

\section{Minimal curved bicomplex on the observer-relative bundle}
\label{sec:bicomplex}

To state the observer-bounded cohomology in which the cacophony
and FRW instances share a cohomology class, we need the
curved bicomplex on the observer-relative symbolic bundle
$(\EE, \hO, \nablaO)$ of the companion manuscript
\texttt{fig2\_companion\_fftc.tex}. Only the minimal structure
required for Def.~\ref{def:ocoh} is given here; the full
bicomplex construction with cocycle placement for each instance
is deferred to post-deadline work.

\begin{definition}[Curved bicomplex on the observer-relative bundle]
\label{def:bicomplex}
Let $\mathcal{U} = \{U_\alpha\}$ be a good cover of the geometric
realization $|\Kt|$, and let $(\EE, \hO, \nablaO)$ be the
observer-relative symbolic bundle over $|\Kt|$ with connection
$\nablaO$ and curvature $\kappaO$ (cf.\ \texttt{fig2\_companion\_fftc.tex},
Def.~1). The \emph{curved bicomplex} is the bigraded space
\[
C^{p, q} \;:=\; \check{C}^p\bigl(\mathcal{U};\, \Omega^q(\cdot, \End(\EE))\bigr),
\]
where $\check{C}^p$ is the \v{C}ech $p$-cochain complex with
values in the sheaf of $\End(\EE)$-valued $q$-forms. The two
differentials are the \v{C}ech coboundary
$d_1 = \delta: C^{p,q} \to C^{p+1, q}$ and the covariant
exterior derivative $d_2 = \DO: C^{p,q} \to C^{p, q+1}$ induced
by $\nablaO$. The total differential is
$D := d_1 + (-1)^p\, d_2$.
\end{definition}

\begin{lemma}[Curvature of the total differential]
\label{lem:curvature}
$D^2 = [\kappaO, \cdot]$ as a map $C^{p,q} \to C^{p, q+2}$. The
bicomplex is \emph{curved}: $D^2 \neq 0$ unless $\kappaO = 0$.
\end{lemma}

\begin{proof}
$d_1^2 = 0$ (standard \v{C}ech), and
$d_2^2 = [\kappaO, \cdot]$ is the curvature of $\nablaO$
acting on $\End(\EE)$-valued forms
(cf.\ \texttt{fig2\_companion\_fftc.tex},
Prop.~curvature-of-nablaO). The cross terms vanish
because $d_1$ acts by restriction and $d_2$ by covariant
differentiation, which commute after sign adjustment; the $(-1)^p$
factor in the definition of $D$ absorbs the sign. Hence
$D^2 = d_1^2 + (-1)^{p+1} d_1 d_2 + (-1)^p d_2 d_1 + d_2^2 =
[\kappaO, \cdot]$.
\end{proof}

\section{The $\OmegaO$-filtration}
\label{sec:filtration}

\begin{definition}[$\OmegaO$-filtration]
\label{def:filtration}
For the curved bicomplex $(C^{\bullet\bullet}, D)$ of
Def.~\ref{def:bicomplex}, the \emph{$\OmegaO$-filtration} is
the decreasing filtration
\[
F^k C^{p, q} \;:=\; \bigl\{\, \omega \in C^{p, q} :
\|\omega(\cdot, \lambda)\|_{\hO} = O(\OmegaO(\lambda)^{k})
\text{ as } \lambda \to \lambda^*\,\bigr\},
\]
where $\lambda$ is the control parameter of an ambient
observer-bounded system and $\omega(\cdot, \lambda)$ denotes the
$\lambda$-parameterized section of $C^{p,q}$. Elements of
$F^k C^{p, q}$ have $\OmegaO$-filtration \emph{order at least}
$k$ (more-vanishing elements have larger $k$).
\end{definition}

\begin{remark}[What the filtration tracks]
$F^k C^{p, q}$ collects cochains whose norm vanishes at least as
fast as $\OmegaO^k$ at the critical point. Since access cost
$\kappaO^{\mathrm{access}}$ diverges as $\OmegaO^{-1/2}$,
elements of $F^k C^{p, q}$ contribute to observables at order
$\OmegaO^{k - 1/2}$ to the observer. Positive-$k$ pieces are
observer-resolvable; negative-$k$ pieces fall below resolution.
\end{remark}

\section{Observer-bounded cohomology}
\label{sec:cohom}

\begin{definition}[Observer-bounded cohomology]
\label{def:ocoh}
The \emph{observer-bounded cohomology} of the $\OmegaO$-filtered
curved bicomplex $(C^{\bullet\bullet}, D, F^\bullet)$ is
\[
\HO^{k, n} \;:=\; H^n\bigl(\mathrm{gr}^k C^{\bullet\bullet},\,
\overline{D}\bigr),
\]
where $\mathrm{gr}^k = F^k / F^{k+1}$ is the associated graded,
and $\overline{D}$ is the induced differential on the graded.
Because $D^2 = [\kappaO, \cdot]$ (Lemma~\ref{lem:curvature})
acts by lowering filtration order (curvature has a definite
$\OmegaO$-order), it descends to zero on $\mathrm{gr}^\bullet$,
and $\overline{D}^2 = 0$. The graded cohomology is therefore
well-defined.
\end{definition}

\begin{principle}[Test-time truth]
\label{princ:tt}
For an observer-bounded framework, the natural cohomology is
$\HO^{\bullet, \bullet}$, not the strict cohomology of
$(C^{\bullet\bullet}, D)$. Strict cohomology requires resolving
differences below the $\OmegaO$-filtration floor, operations
the observer cannot perform. Observer-bounded cohomology
captures \emph{exactly} the content the observer can access.
\end{principle}

\section{Test-time unification}
\label{sec:unification}

\begin{theorem}[Test-time unification]
\label{thm:ttu}
Under observer-bounded cohomology $\HO^{\bullet, \bullet}$:
\begin{enumerate}[label=(\roman*),leftmargin=*,itemsep=2pt]
\item The cacophony obstruction (Prop.~\ref{prop:cac}) admits a
canonical representative $[c_{\Cacph}] \in \HO^{0, 1}$ at
$\OmegaO$-filtration order $0$.
\item The FRW curvature obstruction (Prop.~\ref{prop:frw})
admits a canonical representative $[c_\Frob] \in \HO^{-1, 2}$ at
$\OmegaO$-filtration order $-1$ (sub-resolution at each finite
time, divergent at the critical point).
\item Both classes are functorial images of the same
observer-bounded access functor $\OmegaO$
(Def.~\ref{def:accessor}): they share the same functorial source.
The matter-era $2/3$ arises as the same theorem applied at
$p = 4/3$ instead of $p = 1$; the exponents differ but the
cohomological carrier is the same.
\end{enumerate}
Therefore the cacophony cliff and the FRW past light cone share
a cohomology class in $\HO^{\bullet, \bullet}$ up to
filtration-order shift, with the shift equal to $p$'s difference
from $1$.
\end{theorem}

\begin{proof}[Sketch]
Part (i): the cacophony conflict $c_{ij} = g_i \otimes g_j^* -
g_j \otimes g_i^*$ (constraint pair asymmetry) lives in
$C^{1, 0}$; its $\OmegaO$-filtration order is $0$ because the
cacophony bilinear form vanishes linearly (so $B_\rho$ is
$O(1)$). It is a $d_1$-cocycle (constraint asymmetries are
locally consistent) and its class
$[c_{\Cacph}] \in H^1(\mathrm{gr}^0 C^{\bullet, 0}, \delta)
= \HO^{0, 1}$.
Part (ii): the FRW curvature is
$\kappaO \in \Omega^2(\cdot, \End(\EE))$; its
$\OmegaO$-filtration order is $-1$ (divergent at $t^* = 0$),
and it is the curvature element in $C^{0, 2}$, representing
$[\kappaO] \in \HO^{-1, 2}$.
Part (iii): both are in the image of the access functor
$\OmegaO$ evaluated on their respective systems; the universal
exponent theorem produces the filtration orders from the
vanishing exponents $p$.
\end{proof}

\begin{principle}[ICML/operational band = test-time cohomology band]
\label{princ:icml}
The cohomology $\HO^{\bullet, \bullet}$ is precisely the cohomology
visible to an observer with bounded resolution. Strict cohomology
requires unbounded resolution (hence access below $\OmegaO$'s
minimum), which is not available at inference time. The
operational band of ICML-scale evaluation \emph{is} the test-time
cohomology band.
\end{principle}

\section{The obstruction as the observer bound}
\label{sec:obstruction}

\begin{theorem}[Obstruction = observer bound]
\label{thm:obs-is-bd}
Let $(\EE, \hO, \nablaO)$ be an observer-induced symbolic bundle
(Def.~\ref{def:obs-induced}, below). The following are equivalent:
\begin{enumerate}[label=(\roman*),leftmargin=*,itemsep=2pt]
\item Strict unification closes: there exists a global section
$\alpha \in \Gamma(C^{1, 0} \oplus C^{0, 1})$ with
$D\alpha + \tfrac{1}{2}[\alpha, \alpha] - \kappaO = 0$ in strict
cohomology.
\item The Chern--Weil class $[\kappaO] \in H^2_{dR}(|\Kt|;\,
\End(\EE))$ vanishes.
\item The $\OmegaO$-filtration is trivial: $F^k C^{p, q} = C^{p, q}$
for all $k \in \mathbb{Z}$ (equivalently, $\kappaO \in \bigcap_k F^k C^{0,2}$).
\item Observer resolution is unbounded: $\rmin = 0$.
\end{enumerate}
In particular, the Chern--Weil obstruction to strict decurving
\emph{is} the observer bound, cohomologically expressed. For any
observer-bounded symbolic bundle with nontrivial curvature, (iv)
fails, so (i)-(iii) all fail, and the natural unification is the
test-time unification of Thm.~\ref{thm:ttu}, not the strict one.
\end{theorem}

\subsection*{Preliminary lemmas}

The proof proceeds via a scope assumption on the bundle and
three preliminary results: a filtration-order fact about $\kappaO$
(Lemma~\ref{lem:kappa-filtration}), an \emph{elimination lemma}
showing that both decoupled solvability routes fail under
$\rmin > 0$ (Lemma~\ref{lem:elimination}), and a Maurer--Cartan
solvability statement at $\rmin = 0$
(Lemma~\ref{lem:MC-solvability}).

\begin{definition}[Observer-induced connection]
\label{def:obs-induced}
$\nablaO$ from \texttt{fig2\_companion\_fftc.tex} (Def.~1) is
\emph{observer-induced} if its connection 1-form $\AO$ arises
from the observer's kernel data: $\AO = \AO(\rmin)$ with $\AO \to 0$
in the $\|\cdot\|_{\gO \otimes \hO, \infty}$ norm of
\texttt{fig2\_companion\_fftc.tex} as $\rmin \to 0$. Equivalently,
$\AO$ captures the non-commutation between kernel smoothing and
flat parallel transport at scale $\rmin$; it vanishes in the
Dirac-delta limit $\rmin \to 0$.
\end{definition}

\noindent Lemmas~\ref{lem:kappa-filtration}, \ref{lem:elimination},
\ref{lem:MC-solvability} and Theorem~\ref{thm:obs-is-bd} below
assume an observer-induced connection. Bundles with intrinsic
topological curvature unrelated to $\rmin$ (e.g., pulled back from
a nontrivial classifying map) are outside the scope of
Thm.~\ref{thm:obs-is-bd} as stated; in such cases the Chern--Weil
obstruction persists even at $\rmin = 0$ and (iv) becomes
sufficient but not necessary.

\begin{lemma}[Curvature filtration order]
\label{lem:kappa-filtration}
Let $(\EE, \hO, \nablaO)$ be an observer-induced symbolic bundle
(Def.~\ref{def:obs-induced}) with kernel scale $\rmin$ and
$\kappaO = d\AO + \AO \wedge \AO$. Under the $\OmegaO$-filtration
of Def.~\ref{def:filtration}:
\begin{enumerate}[label=(\alph*),leftmargin=*,itemsep=2pt]
\item $\kappaO \in F^0 C^{0, 2}$: the curvature norm is bounded
uniformly in $\lambda$.
\item $\kappaO \in \bigcap_{k \in \mathbb{Z}} F^k C^{0, 2}$ if and
only if $\kappaO = 0$ identically.
\item Under the observer-induced hypothesis, $\kappaO = 0$
identically if and only if $\rmin = 0$.
\end{enumerate}
\end{lemma}

\begin{proof}
(a) $\|\kappaO\|_{\hO} \leq \|d\AO\|_{\hO}
+ \|\AO\|^2_{\gO \otimes \hO, \infty}$. Both terms are bounded
uniformly in $\lambda$: the second by $\Obs$-boundedness of $\AO$
(\texttt{fig2\_companion\_fftc.tex}, Def.~1), the first by
smoothness of $\AO$ on the open patch
$|\Kt| \subset (M, g)$. Hence $\|\kappaO(\cdot, \lambda)\|_{\hO}
= O(1) = O(\OmegaO(\lambda)^0)$.

(b) ($\Leftarrow$) The zero cochain lies in every $F^k$ trivially.
($\Rightarrow$) If $\kappaO \in F^k$ for all $k$, then
$\|\kappaO(\cdot, \lambda)\|_{\hO} = O(\OmegaO(\lambda)^k)$ for
every $k > 0$; since $\OmegaO(\lambda) \to 0$ at $\lambda^*$, the
norm decays faster than every polynomial order. By continuity of
$\kappaO$ in $\lambda$ (inherited from smoothness of $\AO$),
$\kappaO \equiv 0$ at $\lambda^*$; smoothness in $\lambda$ then
propagates the identity in a neighborhood, and real-analyticity
of the construction extends it globally.

(c) ($\Leftarrow$) If $\rmin = 0$, $\AO \to 0$ by
Def.~\ref{def:obs-induced}, so
$\kappaO = d\AO + \AO \wedge \AO \to 0$ in the norm of~(a);
the limit object is the zero 2-form on $|\Kt|$.
($\Rightarrow$) If $\kappaO = 0$ and $\AO$ is observer-induced,
then $\AO$ is flat for every $\rmin$ in its parameter range.
For observer-induced $\AO$, flatness at all $\rmin$ forces
$\AO \equiv 0$ (the kernel-induced non-commutation is nonzero
for $\rmin > 0$ by the construction in
\texttt{fig2\_companion\_fftc.tex}); which in turn forces
$\rmin = 0$ via Def.~\ref{def:obs-induced}.
\end{proof}

\begin{lemma}[Elimination: no strict-complex MC solvability strategy survives under $\rmin > 0$]
\label{lem:elimination}
Suppose $\rmin > 0$, $\AO$ is observer-induced
(Def.~\ref{def:obs-induced}), and the good cover
$\mathcal{U} = \{U_\beta\}$ has chart radius commensurable with
the observer kernel scale ($\mathrm{diam}_g(U_\beta) \asymp \rmin$).
Consider candidate Maurer--Cartan forms
$\alpha = \alpha_{\mathrm{C}} + \alpha_{\mathrm{R}}
\in C^{1,0} \oplus C^{0,1}$ solving
$D\alpha + \tfrac{1}{2}[\alpha, \alpha] = \kappaO$
in the \emph{strict} bicomplex $(C^{\bullet, \bullet}, D)$. The
two natural proof strategies --- local-first (R1) and
global-first (R2) --- both fail at the Chern--Weil class
$[\kappaO]$:
\begin{enumerate}[label=(R\arabic*),leftmargin=*,itemsep=2pt]
\item \textbf{Local-first (de Rham-first).} Construct local
covariant primitives
$\tilde\alpha_{\mathrm{R}}|_{U_\beta}$ on each chart, then patch.
Local primitives \emph{exist} on each contractible
kernel-scale chart (by covariant integration), but they fail to
glue: their \v{C}ech descent defect
$\beta := \delta \tilde\alpha_{\mathrm{R}} \in C^{1,1}$ represents,
under the \v{C}ech--de Rham isomorphism, the class
$[\kappaO] \in H^2_{dR}(|\Kt|;\, \End(\EE))$, which is nonzero
under the hypotheses. No global
$\alpha_{\mathrm{C}} \in C^{1,0}$ satisfies $\delta \alpha_{\mathrm{C}}
= -\beta$.
\item \textbf{Global-first (\v{C}ech-first).} Trivialize the
bundle globally via a \v{C}ech cocycle (gauge transformation
across charts) before integrating. Fails at the initial step:
the Chern--Weil class $[\kappaO] \neq 0$ is exactly the obstruction
to global flat trivialization. No global
$\alpha_{\mathrm{C}}$ reduces the connection to a flat
representative.
\end{enumerate}
R1 and R2 are \emph{two presentations of the same underlying
class obstruction} $[\kappaO] \neq 0$, linked by \v{C}ech--de Rham;
they are not logically independent. The content of the lemma is
that \emph{every proof strategy in the strict bicomplex encounters
this class at some step} --- local-first at gluing, global-first
at trivialization. Neither ordering yields a strict-complex
solution; the coupled $\OmegaO$-filtered framework is the natural
survivor.
\end{lemma}

\begin{proof}
\textbf{(R1) failure --- local-first fails at gluing.}
On each contractible kernel-scale chart $U_\beta$, the ambient
smoothness of $\AO$ and $\kappaO$ gives a local
$\tilde\alpha_{\mathrm{R}}|_{U_\beta} \in \Omega^1(U_\beta;\, \End(\EE))$
solving $\DO \tilde\alpha_{\mathrm{R}} = \kappaO|_{U_\beta}$ via
covariant integration along radial paths from a basepoint
(explicit in coordinates: the radial-gauge primitive). Passing
to the cover $\mathcal{U}$, the collection
$\{\tilde\alpha_{\mathrm{R}}|_{U_\beta}\}_\beta$ has pairwise
differences on overlaps giving a \v{C}ech 1-cocycle
$\beta := \delta \tilde\alpha_{\mathrm{R}} \in C^{1,1}$. Under
the \v{C}ech--de Rham isomorphism
(cf.\ Bott--Tu, Ch.~8; Kobayashi--Nomizu, Ch.~III.5), the class
of $\beta$ in $\check{H}^1(\mathcal{U};\, \Omega^1(\End(\EE)))$
corresponds to $[\kappaO] \in H^2_{dR}(|\Kt|;\, \End(\EE))$.
Nonzero class (established below) means $\beta$ is not a
\v{C}ech coboundary, so no $\alpha_{\mathrm{C}}$ absorbs it.

\textbf{(R2) failure --- global-first fails at trivialization.}
Global trivialization of $\nablaO$ to a flat connection via a
\v{C}ech cocycle requires $[\kappaO] = 0$ in
$H^2_{dR}(|\Kt|;\, \End(\EE))$ by the standard Chern--Weil
obstruction. Under nonzero class, no such trivialization exists.

\textbf{$[\kappaO] \neq 0$ under observer-induced $\rmin > 0$.}
Contrapositive of Lemma~\ref{lem:kappa-filtration}(c): if
$[\kappaO] = 0$, Chern--Weil supplies a global flat gauge, forcing
$\AO \equiv 0$ (in some gauge), which by the observer-induced
hypothesis (Def.~\ref{def:obs-induced}) forces $\rmin = 0$.
Contrary to assumption.

\textbf{Relation between the two failures.} R1 and R2 encounter
the class $[\kappaO]$ at different proof steps --- R1 after
local primitives are constructed, R2 before any construction is
attempted --- but the class itself is a single Čech--de Rham
invariant. No proof strategy in the strict bicomplex avoids it,
which is the elimination content.
\end{proof}

\begin{remark}[Proof by elimination and finite capture]
\label{rem:elim-structural}
Lemma~\ref{lem:elimination} has the structural shape of
dual-horizon arguments in observer-bounded epistemology
(cf.\ \texttt{main.tex}, \S on dual horizons): two
single-direction proof strategies meet the same obstruction at
distinct steps, leaving the coupled (filtered) framework as the
only surviving solvability path. The two failure points --- R1
\emph{local differential}, R2 \emph{global topological} --- are
dual presentations of one underlying class obstruction, but they
block two distinct proof flows. The creative content of the
lemma is not that R1 and R2 are logically independent (they
aren't; \v{C}ech--de Rham ties them), but that \emph{no
single-direction proof strategy evades the obstruction}, and
finite capture ($\rmin > 0$) is what puts the obstruction class
in place to begin with: on the observer-induced kernel-scale
cover, the Chern--Weil class is generically nonzero precisely
because the observer resolves only at scale $\rmin$. We record
this as a structural pattern: \emph{observer finite capture
creates the cohomology that the coupled framework is forced to
carry}.
\end{remark}

\begin{lemma}[Observer-bounded Maurer--Cartan solvability]
\label{lem:MC-solvability}
If $\rmin = 0$, the Maurer--Cartan equation
$D\alpha + \tfrac{1}{2}[\alpha, \alpha] - \kappaO = 0$ admits a
global solution
$\alpha \in \Gamma(C^{1, 0} \oplus C^{0, 1})$.
\end{lemma}

\begin{proof}
By Lemma~\ref{lem:kappa-filtration}(c), $\rmin = 0$ gives
$\kappaO = 0$ identically. The equation reduces to
$D\alpha + \tfrac{1}{2}[\alpha, \alpha] = 0$, solved trivially by
$\alpha = 0$. (Equivalently, the observer-relative bundle at
$\rmin = 0$ is flat, and any flat connection is gauge-equivalent
to the trivial one on a simply connected base.)
\end{proof}

\subsection*{Proof of Theorem~\ref{thm:obs-is-bd}}

\begin{proof}
We establish the cycle
(iv) $\Rightarrow$ (iii) $\Rightarrow$ (ii) $\Rightarrow$ (i)
$\Rightarrow$ (iv).

\textbf{(iv) $\Rightarrow$ (iii).}
Assume $\rmin = 0$. By Lemma~\ref{lem:kappa-filtration}(c),
$\kappaO = 0$, and by the observer-induced hypothesis $\AO = 0$
as well (Def.~\ref{def:obs-induced}). The flat trivial
observer-relative connection makes every cochain
$\omega \in C^{p,q}$ $\lambda$-independent at leading order
(no $\OmegaO$-coupling arises), so
$\|\omega(\cdot, \lambda)\|_{\hO} = O(\OmegaO(\lambda)^k)$ for
every $k \in \mathbb{Z}$. Hence $F^k C^{p, q} = C^{p, q}$ for
all $k$; in particular, $\kappaO \in \bigcap_k F^k$.

\textbf{(iii) $\Rightarrow$ (ii).}
Assume $\kappaO \in \bigcap_k F^k C^{0, 2}$. By
Lemma~\ref{lem:kappa-filtration}(b), $\kappaO = 0$ identically.
A vanishing curvature 2-form has trivial de Rham class:
$[\kappaO] = 0$ in $H^2_{dR}(|\Kt|;\, \End(\EE))$.

\textbf{(ii) $\Rightarrow$ (i).}
Assume $[\kappaO] = 0$. By the standard Chern--Weil construction
(cf.\ Kobayashi--Nomizu, Bott--Tu; applied here to the
observer-relative bundle as in \texttt{fig2\_companion\_fftc.tex}),
trivial curvature class implies the existence of a connection
1-form shift $\alpha$ whose action $\nablaO \to \nablaO + \alpha$
produces a flat connection. This is equivalent to solving
$\DO \alpha + \tfrac{1}{2}[\alpha, \alpha] = \kappaO$ locally.
Globalization of $\alpha$ across the good cover $\mathcal{U}$
uses the vanishing of the \v{C}ech obstruction (also guaranteed by
$[\kappaO] = 0$) together with a partition of unity, yielding
$\alpha \in \Gamma(C^{1,0} \oplus C^{0,1})$ satisfying
$D\alpha + \tfrac{1}{2}[\alpha, \alpha] - \kappaO = 0$ in the
strict bicomplex.

\textbf{(i) $\Rightarrow$ (iv).}
By contrapositive: assume $\rmin > 0$, and show (i) fails. By
Lemma~\ref{lem:elimination}, the strict bicomplex admits no
global MC solution when $\rmin > 0$: the de Rham-first route R1
fails at local covariant integration, the \v{C}ech-first route
R2 fails at global topological descent, and no combined decoupling
in the strict complex evades the class obstruction
$[\kappaO] \neq 0$ that both present. Any MC solvability content
that survives lives in the $\OmegaO$-filtered (coupled)
framework, absorbing $\kappaO$ at nontrivial filtration order ---
the observer-bounded cohomology content of Thm.~\ref{thm:ttu},
not a strict global section in
$\Gamma(C^{1,0} \oplus C^{0,1})$. Hence (i) fails, and the
contrapositive gives (i) $\Rightarrow$ $\rmin = 0 =$~(iv).
\end{proof}

\begin{remark}[Why the cycle goes through elimination]
\label{rem:why-elim}
The implication (i) $\Rightarrow$ (iv) can formally be derived
through (ii) and (iii) without invoking
Lemma~\ref{lem:elimination}. We have routed through the
elimination lemma because it gives the \emph{direct} obstruction:
strict MC solvability fails under $\rmin > 0$ not because of a
topological accident, but because neither single-direction route
(de Rham-first, \v{C}ech-first) can absorb $\kappaO$. The coupling
between observer resolution and coupled cohomology is therefore
not a byproduct of the Chern--Weil isomorphism --- it is a
structural consequence of finite capture.
\end{remark}

\begin{remark}[Thm.~\ref{thm:obs-is-bd} as a screening statistic]
\label{rem:screening}
Thm.~\ref{thm:obs-is-bd} is formulated as a
\emph{screening statistic}, matching the falsifiability posture
of the companion ICML paper~\cite{tiffany2026cacophony} and
§\ref{sec:falsifiability} below. Rather than a blanket
``obstruction = observer bound'' across all symbolic bundles,
it discriminates two regimes:
\begin{enumerate}[label=(\alph*),leftmargin=*,itemsep=2pt]
\item \emph{Observer-induced} (Def.~\ref{def:obs-induced}): the
Chern--Weil class exists \emph{because of} finite capture. Here
obstruction $\equiv$ observer bound, and $\rmin \to 0$ dissolves
both.
\item \emph{Intrinsic-topological}: the class exists
independently of $\rmin$ (e.g., pulled back from a nontrivial
classifying map). Here the observer bound is orthogonal to the
obstruction; $\rmin \to 0$ leaves $[\kappaO] \neq 0$.
\end{enumerate}

\smallskip\noindent\emph{Operational diagnostic.}
Given a system with measured $[\kappaO] \neq 0$, vary $\rmin$
monotonically. If $[\kappaO] \to 0$ as $\rmin \to 0$, the bundle
is observer-induced and the theorem applies. If $[\kappaO]$
stays bounded away from zero as $\rmin \to 0$, the bundle carries
intrinsic topology and the observer bound is not the carrier of
the class. The theorem's content is this discrimination; the
screening is intentional specificity, not a hedge. The
cacophony and FRW instances of \S\ref{sec:cacophony} and
\S\ref{sec:frw} are observer-induced by construction: both have
topologically trivial bases, so $[\kappaO] = 0$ away from the
critical point, and the obstruction that Thm.~\ref{thm:ttu}
packages into $\HO^{\bullet, \bullet}$ is precisely the
finite-capture class.
\end{remark}

\section{Falsification criteria}
\label{sec:falsifiability}

The Universal Exponent Theorem (Thm.~\ref{thm:universal}) and
its two instances are falsifiable:

\textbf{F1} (cacophony). If, in the cacophony
bilinear family, some rank-one deformation produces a critical
exponent $\neq -1/2$ (i.e., vanishing order $p \neq 1$), the
cacophony instance is misclassified. Testable via variant
constraint geometries.

\textbf{F2} (cosmology). If a perfect-fluid FRW cosmology with
known $w$ yields a measured access-cost exponent differing from
$-2/(3(1+w))$ (Cor.~\ref{cor:genw}), the FRW instance is
misclassified or the functor is not $\OmegaO$. Testable against
Planck Euclid BAO joint constraints.

\textbf{F3} (universality). If there is any observer-bounded
system whose capacity vanishes to order $p$ and whose
access-cost exponent is not $-p/2$, Thm.~\ref{thm:universal} is
wrong.

\textbf{F4} (cohomological unification). If strict cohomological
lift is constructed for the cacophony--FRW pair without
requiring $\rmin = 0$, Thm.~\ref{thm:obs-is-bd} is wrong.

\section{Revisions to main-text language}
\label{sec:revisions}

The main-text changes supported by this companion are two:

\textbf{Caption (Fig.~3).} The earlier caption phrased radiation-era
$1/2$ as a match and matter-era $2/3$ as regime-distinct. Under
Thm.~\ref{thm:universal} both are instances of one theorem with
different $p$; the revised caption reports the prediction
$-2/(3(1+w))$ and attributes the $1/2$ match to $w = 1/3$.

\textbf{H$^1$ claim (Empirical Concurrence, \texttt{main.tex} line 1098).}
The earlier sentence suggested $H^1$ \emph{may be} a shared
invariant. Under Thm.~\ref{thm:ttu} the shared invariant lives
in observer-bounded cohomology (Def.~\ref{def:ocoh}), and under
Thm.~\ref{thm:obs-is-bd} the failure of a strict $H^1$ lift is
equivalent to the observer bound. The revised sentence states
this precisely, citing Thm.~\ref{thm:universal} and
Def.~\ref{def:ocoh} of this companion.

\section*{Bibliography (abbreviated)}
\label{sec:bib}

\begin{thebibliography}{9}

\bibitem{tiffany2026cacophony}
Anonymous.
\emph{The Cost of Cacophony: Geometric Limits on Multi-Constraint
Alignment}. Proc. ICML, 2026.

\bibitem{planck2020parameters}
Planck Collaboration.
\emph{Planck 2018 results. VI. Cosmological parameters}. A\&A 641
(2020) A6.

\bibitem{tiffany2026hypsurf}
Anonymous.
\emph{The Hypothesis Surface: An Operational Epistemology for
Autonomous Research}. AGI-26, 2026 (companion main manuscript).

\bibitem{fig2comp}
Anonymous.
\emph{Companion Manuscript, Figure~2: FFTC and Observer-Relative
Stokes}. File \texttt{fig2\_companion\_fftc.tex}.

\end{thebibliography}

\end{document}
